\documentclass[12pt]{article}
\usepackage[utf8]{inputenc}
\usepackage[T1]{fontenc}
\usepackage{hyperref}
\title{Xiri\^ana Language Resources for K-12 Curriculum Development}
\author{Compiled by Claude Sonnet 5.5 (claude-sonnet-5-5) for Andr\'as Kornai}
\date{}
\begin{document}
\maketitle

\section{Introduction}
Xiri\^ana (ISO 639-3 xir; Glottocode xiri1243) is an Arawakan language of Amazonas state, Brazil. It is also called Shiriana, Chiriana, Ciriana or Bahuana. Search results describe it as extinct or dormant, with a date of about 2000 given in connection with the death of Claudina de Souza Soares. Sources disagree on its classification. One places it in Central Arawakan and most closely related to Manao and Kariai. Another calls it an unclassified Upper Amazon Arawakan language. Glottolog files it under Negro-Roraima and Bahuanaic. Features reported include no grammatical gender and an active--stative syntax. The search found no school materials, primers, or revitalization programme. The search also found no information on the current status of the Shiriana communities.

\section{Primary Educational Resources}
The search found no primers, textbooks or school programmes for Xiri\^ana. Any curriculum would have to be built from the scholarly documentation listed below.

\section{Linguistic and Reference Materials}
\begin{itemize}
\item Wikipedia, ``Xiriana language'': \url{https://en.wikipedia.org/wiki/Xiriana_language}
\item Wikipedia, ``Bahuana language'': \url{https://en.wikipedia.org/wiki/Bahuana_language}
\item Wikipedia, ``Shiriana'' (the people): \url{https://en.wikipedia.org/wiki/Shiriana}
\item Glottolog entry for xiri1243: \url{https://glottolog.org/resource/languoid/id/xiri1243.bigmap.html}
\item The main scholarly reference named in the results is Henri Ramirez (1992), ``Le Bahuana: une nouvelle langue de la famille arawak'', \emph{Amerindia}. The search showed no URL for it.
\end{itemize}

\section{Religious and Community Resources}
The search found none.

\section{Digital and Audio Resources}
The search found no audio, video or corpus material.

\section{Curriculum Development Context}
Very little is documented. A heritage curriculum would have to be a reclamation project based on Ramirez's description and any other archival data, and would need the involvement of the Shiriana community. This should be treated as exploratory work, not as a comprehensive curriculum.
\end{document}